% Shared visual style for diagrams and algorithms.
% Included once from main.tex; not a chapter, contains no \chapter command.

\usepackage{tikz}
\usetikzlibrary{arrows.meta,positioning,calc,decorations.pathreplacing,patterns,matrix,fit,backgrounds,shapes.geometric}
\usepackage{pgfplots}
\pgfplotsset{compat=1.18}
\usepackage{float}
\usepackage{subcaption}
\usepackage{algorithm}
\usepackage{algpseudocode}

% ---- Palette: restrained, print-friendly, colorblind-conscious ----
\definecolor{pqcblue}{HTML}{1F4E79}   % primary structure, arrows, key vectors
\definecolor{pqcteal}{HTML}{2E7D6B}   % secondary structure
\definecolor{pqcorange}{HTML}{C1601B} % emphasis / the "answer" or highlighted object
\definecolor{pqcgray}{HTML}{6B6B6B}   % grid lines, muted context
\definecolor{pqclight}{HTML}{EDEDED}  % fills / backgrounds

% ---- Interactive protocol figures ----
% The convention of Lyubashevsky, Nguyen and Plancon (CRYPTO 2022), which is
% the house style of the lattice zero-knowledge literature: one framed figure
% carrying a header of public and private information and the relation being
% proved, a rule, then Prover and Verifier columns with labelled message
% arrows between them and an "Accept iff" block closing the verifier's side.
% Preferred over numbered pseudocode whenever who-knows-what and who-sends-when
% are part of what the reader has to see.
\newsavebox{\pqcprotobox}
\newenvironment{protocol}
  {\begin{lrbox}{\pqcprotobox}\begin{minipage}{0.93\textwidth}\small\vspace{2pt}}
  {\vspace{2pt}\end{minipage}\end{lrbox}\centering
   {\setlength{\fboxsep}{6pt}\fbox{\usebox{\pqcprotobox}}}}
% Labelled message arrows. The arrow stretches to the width of whatever
% column it sits in, so a send visually spans the gap between the two
% parties instead of floating as a fixed-length stub; the payload rides
% above the shaft, as in the LNP figures.
\newcommand{\msgto}[1]{\ensuremath{\xrightarrow{\makebox[0.88\linewidth][c]{$\scriptstyle #1$}}}}
\newcommand{\msgfrom}[1]{\ensuremath{\xleftarrow{\makebox[0.88\linewidth][c]{$\scriptstyle #1$}}}}
% The header/interaction separator, and the column heads.
\newcommand{\protorule}{\par\smallskip\hrule\smallskip}
\newcommand{\Prover}{\underline{Prover}}
\newcommand{\Verifier}{\underline{Verifier}}

% ---- Two-party message-sequence diagrams ----
% One vocabulary for every "A sends, B replies" picture in the book (the
% KEM handshake, Fiat--Shamir, the TLS hybrid): bold party names over
% dashed lifelines, protocol messages as blue arrows whose emission point
% is marked on the sending lifeline, local computation as muted side
% notes, and an orange result strip when the parties end up sharing
% something. Blue is reserved for what actually crosses the wire, so a
% reader can tell a message from a structural arrow at a glance.
\tikzset{
  seqparty/.style={font=\small\bfseries},
  seqlife/.style={pqcgray!45, dashed},
  seqmsg/.style={{Circle[length=3pt]}-{Latex[length=2.5mm]}, pqcblue, semithick},
  seqnote/.style={pqcgray},
  seqresult/.style={keynode, align=center},
}

% ---- Reusable TikZ styles ----
\tikzset{
  latticept/.style={circle, fill=pqcgray, inner sep=0pt, minimum size=3pt},
  latticeptbold/.style={circle, fill=pqcblue, inner sep=0pt, minimum size=4pt},
  basisvec/.style={-{Latex[length=2.5mm]}, pqcblue, thick},
  altvec/.style={-{Latex[length=2.5mm]}, pqcteal, thick},
  shortvec/.style={-{Latex[length=2.5mm]}, pqcorange, thick, line width=1pt},
  flownode/.style={draw=pqcblue, thick, rounded corners=2pt, fill=pqclight,
                    minimum height=7mm, inner sep=3pt, align=center, font=\small},
  flowarrow/.style={-{Latex[length=2.5mm]}, pqcgray, thick},
  keynode/.style={draw=pqcorange, thick, rounded corners=2pt, fill=white,
                   minimum height=7mm, inner sep=3pt, align=center, font=\small},
}

% ---- Algorithm captions read "Algorithm N: Name" consistently ----
\floatname{algorithm}{Algorithm}
\renewcommand{\algorithmicrequire}{\textbf{Input:}}
\renewcommand{\algorithmicensure}{\textbf{Output:}}

% ---- Code and diagrams never straddle a page break ----
% TikZ pictures, tabulars and [H] floats are unbreakable boxes already;
% a code listing is the one display element TeX will happily split across
% pages. Boxing every listing in a minipage forbids the split, and with
% \raggedbottom below, a listing that no longer fits simply moves whole
% to the next page. Safe because every listing in the book is well under
% a page (the longest is 43 lines); a listing taller than \textheight
% would overflow the bottom margin and must be split up at the source.
\usepackage{etoolbox}
\BeforeBeginEnvironment{lstlisting}{\par\noindent\begin{minipage}{\linewidth}}
\AfterEndEnvironment{lstlisting}{\end{minipage}\par}

% ---- Vertical spacing ----
% book.cls flush-bottoms every page by stretching whatever glue it offers,
% which on pages full of [H] floats and short paragraphs inflates the display
% and list skips unevenly -- the same gap is a different height on every page.
\raggedbottom
% The chapters alternate one-sentence paragraphs with one-line display
% formulas; at the default skips each formula reads as a paragraph of its own
% inside a blank line on either side. Tighten the skips so a short formula
% stays visually attached to the sentence that introduces it. Appended to
% \normalsize because every size command resets these lengths.
\makeatletter
\g@addto@macro\normalsize{%
  \setlength{\abovedisplayskip}{6pt plus 2pt minus 3pt}%
  \setlength{\belowdisplayskip}{6pt plus 2pt minus 3pt}%
  \setlength{\abovedisplayshortskip}{2pt plus 1pt}%
  \setlength{\belowdisplayshortskip}{4pt plus 2pt minus 2pt}%
}
\makeatother
% Lists inherit the same airiness from the class defaults.
\setlist{topsep=5pt, itemsep=2pt, parsep=0pt}

% ---- Theorem-like environments (shared counter, numbered by chapter) ----
\usepackage{amsthm}
\theoremstyle{definition}
\newtheorem{definition}{Definition}[chapter]
\newtheorem{example}[definition]{Example}
\theoremstyle{plain}
\newtheorem{theorem}[definition]{Theorem}
\newtheorem{proposition}[definition]{Proposition}
\newtheorem{lemma}[definition]{Lemma}
\theoremstyle{remark}
\newtheorem*{remark}{Remark}

% ---- Running headers show "Chapter N. Title" on every page, not section
% titles. Long \section{} titles otherwise collide with the page number
% when they land on a recto page; showing chapter context on both sides
% is standard for a book this size and sidesteps that failure mode
% entirely. Individual chapters can still shorten their own header with
% \chaptermark{...} right after \chapter{...} (see e.g. 14_cbd.tex).
\makeatletter
\renewcommand{\sectionmark}[1]{}
\renewcommand{\chaptermark}[1]{%
  \markboth{\MakeUppercase{\ifnum \c@secnumdepth >\m@ne \@chapapp\ \thechapter.\ \fi #1}}%
           {\MakeUppercase{\ifnum \c@secnumdepth >\m@ne \@chapapp\ \thechapter.\ \fi #1}}%
}
\makeatother

% Reference to a companion Sage file, hyperlinked to the repository.
\newcommand{\sagerepo}{https://github.com/AppliedPQC/AppliedPQC/blob/main/sage}
\newcommand{\sagefile}[1]{\href{\sagerepo/#1}{\texttt{sage/#1}}}

% Pointer to the browser-runnable version of the companion Sage code.
\newcommand{\playgroundurl}{https://appliedpqc.io/playground.html}
\newcommand{\playground}[1]{%
  \smallskip\noindent\emph{Run it in a browser:} #1. See
  \href{\playgroundurl}{appliedpqc.io/playground.html} --- nothing to install.}

